% Experiment 07 -- the PKIX-underuse longitudinal panel (PANEL): standalone writeup.
% IEEEtran conference format, matching paper/paper.tex conventions.
% Compile: pdflatex exp07_panel.tex   (figures, once generated, live in paper/figures/).
%
% STATUS: data collection is LIVE (launched 2026-07-11, auto-stops 2026-07-18). The Methodology
% is complete; the Analysis section (Sec. V) is deliberately SCAFFOLDED -- each subsection states
% what it will report and carries a \pending marker and a commented figure slot, to be filled once
% the 7-day window closes and results/a (traceroutes) + results/b (pings) are merged.
\documentclass[conference]{IEEEtran}
\IEEEoverridecommandlockouts
\usepackage{cite}
\usepackage{graphicx}
\usepackage{booktabs}
\usepackage{amsmath,amssymb}
\usepackage{xcolor}
\usepackage[hidelinks]{hyperref}
\usepackage{url}
\graphicspath{{figures/}}

% --- markers for results that land only after the 7-day run completes ---
\newcommand{\pending}{{\color{gray}\small[pending: run completes 2026-07-18]}}
\newcommand{\XX}{{\color{gray}$\langle$XX$\rangle$}} % numeric placeholder to fill from data

\begin{document}

\title{A Uniform Longitudinal Panel of Domestic and\\
Offshore Routing Quality in Pakistan}

\author{\IEEEauthorblockN{Rayan Atif}
\IEEEauthorblockA{Lahore University of Management Sciences (LUMS)\\
Supervised by Dr.\ Saqib Ilyas~~$\cdot$~~Funded by the APNIC Foundation\\
\texttt{rayan.atif2@gmail.com}}}

\maketitle

\begin{abstract}
Our earlier measurements established \emph{that} Pakistani domestic traffic frequently leaves the
country and that offshore hosting is common, but each was a single snapshot taken with a different
set of probes, a different traceroute protocol, and a different definition of round-trip time, so
none of them can support a fair, per-ISP comparison of the user experience or say whether the
penalty is a fixed structural cost or a passing artefact of the hour it was measured. This paper
describes \textbf{PANEL}, a longitudinal measurement built to remove exactly those inconsistencies:
from \emph{every} connected Pakistani RIPE Atlas probe, to a \emph{fixed, reproducible} stratified
sample of $100$ websites, we run one TCP/80 Paris traceroute per hour and one ping every
$30$~minutes continuously for seven days, under a single framework that reports the same four
quality indicators---RTT, hop count, jitter, and packet loss---for every vantage. The same-method,
same-target, many-days design is what lets us (i)~compare hosting classes and ISPs on equal terms,
(ii)~measure how \emph{stable} a tromboning verdict is over hours and days rather than trusting one
trace, (iii)~separate any time-of-day or weekday cycle from a genuinely structural penalty, and
(iv)~hold a normal-day baseline ready so that the next submarine-cable fault is captured with a
true before/during/after. We report the design, the sampling rationale, and the operational
scheme that fits RIPE Atlas's per-account concurrency limit; the seven-day dataset is in
collection at the time of writing, and the per-ISP distributions and stability analyses it feeds
are presented in Section~\ref{sec:analysis}. \pending
\end{abstract}

\section{Introduction}
Whether Pakistan's Internet Exchange (PKIX) is worth using is ultimately a question about the
\emph{user}: do the customers of an ISP that peers locally get measurably better, and more stable,
service than the customers of one that does not? Answering it requires more than establishing that
hairpinning exists---which our prior work did, finding roughly $40\%$ of top Pakistani sites hosted
offshore and about $11\%$ of small-ISP address space reachable only by detouring abroad. It
requires a comparison that is \emph{fair}: the same targets, measured the same way, from every ISP,
often enough to distinguish a structural cost from a momentary one. Our earlier experiments cannot
provide that comparison, and the reason is instructive. The vantage set changed from one run to the
next (five probes, then three, then eight, then one, then fourteen); the path protocol changed from
ICMP to TCP/80; RTT was variously a single packet, an average, or a maximum hop; and only the
outage study carried jitter and loss at all. Each choice was locally reasonable, but together they
make cross-experiment numbers incommensurable and every result a single snapshot, fragile to the
intermittency we ourselves observed when a verdict flipped minutes apart.

This paper closes that gap with a deliberately uniform, longitudinal design. Rather than add another
snapshot, \textbf{PANEL} fixes the measurement---one probe set (all connected Pakistani probes,
discovered live), one protocol (TCP/80 Paris traceroute for the path, ping for RTT), one RTT
definition (the minimum of $N$ packets), and one four-indicator metric set---and then holds it
constant for seven continuous days against a frozen, reproducible list of $100$ websites. Holding
everything constant is precisely what turns previously incomparable observations into a coherent
measurement: because only the vantage and the target vary, a difference in the data is a difference
in the network, not in the instrument. Concretely, PANEL is built to deliver four things the
snapshots could not: a per-ISP distribution of each quality indicator to local, CDN, and offshore
targets (RQ1); a measurement of \emph{same-ISP} consistency across the several probes we hold per
operator (RQ2); a \emph{stability} profile for tromboning, quantifying how often a given path's
verdict changes over the week; and a standing normal-day baseline that lets a future cable fault be
read as a true departure from normal rather than, as before, against an already-degraded state. The
contribution of this paper is the design and its justification together with the operational scheme
that makes it runnable under RIPE Atlas's limits; the resulting dataset is in collection, and the
analyses it feeds are laid out, with their figures reserved, in Section~\ref{sec:analysis}.

\section{Related Work}
The study most like ours in shape is Di Bartolomeo et al.~\cite{dibartolomeo2015}, who use national
RIPE Atlas probes to compare IXP against upstream paths for Italian ISPs and report the same four
indicators---RTT, hop count, jitter, packet loss---as cumulative distributions, one curve with the
exchange and one without. We adopt both their metric set and their figure language, but where they
compare two BGP configurations of the same probes at one time, we compare \emph{many probes of
many ISPs} to the same targets \emph{over many days}, trading their controlled announcement for the
statistical leverage of a longitudinal panel. That the underlying phenomenon is real beyond Italy is
clear from Gupta et al.~\cite{gupta2014}, who show African ISPs that do not peer locally sending
domestic traffic on a detour through Europe---the same hairpinning we quantify in Pakistan, and the
motivation for measuring its cost rather than merely its presence. The value of measuring
\emph{over time} rather than once follows the longitudinal tradition in inter-domain
measurement~\cite{sanchez2014}, in which a repeated observation exposes structure---diurnal cycles,
stability, rare events---that no single snapshot can. Our platform and resolution tooling are the
community standard: RIPE Atlas for the vantages~\cite{ripeatlas}, RIPEstat~\cite{ripestat} and Team
Cymru with an RDAP fallback for AS and registry attribution, and the sampling of announced prefixes
in the spirit of selective address-space scanning~\cite{klick2016}. Large-IXP
studies~\cite{ager2012,chatzis2013} supply the background on why local peering matters at all.

\section{Methodology}
\label{sec:method}
PANEL is one measurement framework applied uniformly; this section states each of its choices and,
because the design \emph{is} the contribution at this stage, why each is what it is.

\subsection{Vantages: every connected Pakistani probe}
We source probes not from a hardcoded list but from a live query to the RIPE Atlas API for all
connected probes in Pakistan (\texttt{country\_code=PK, status=1}) at schedule time---seventeen at
launch. Discovering the set live is what keeps the panel comparable to itself as probes come and go,
and it maximises coverage on the two axes the research questions need: \emph{across} ISPs, so that
Set-3 (locally peering) operators can be compared with Set-1/2 operators, and \emph{within} an ISP,
since we hold several probes for the same operator (Cybernet~$\times3$, PTCL~$\times3$,
Nayatel~$\times2$) which is exactly what RQ2's same-ISP consistency test consumes. Three probes are
known to distort path metrics---two are ICMP-filtered and one runs in a container that exposes only
two hops---so their hop counts are excluded from path-length statistics while their RTT, which
remains valid, is retained; we carry these exclusions explicitly rather than silently dropping the
probes.

\subsection{Targets: a two-axis stratified sample of 100 sites}
\label{sec:targets}
The panel measures websites only---no canaries and no bare IP blocks---because its question is the
quality a user actually experiences loading real content, and dedicated block-level tromboning
intermittency is the job of a separate experiment. The target set is a \emph{frozen} sample of
$100$ sites (\texttt{targets.csv}), drawn so that it is representative on one axis and powerful on
the other. On the first axis, \textbf{genre}, we stratify by CISA critical-infrastructure sector
\emph{proportionally} to the sector mix of the resolved \texttt{.pk} web population, so that a
pooled or reweighted statistic generalises to ``the Pakistani web'' without post-hoc correction
(Table~\ref{tab:sectors}). On the second axis, \textbf{hosting class}, we deliberately depart from
the population and impose a fixed \textbf{$40$ Pakistani / $40$ CDN / $20$ Abroad} split, held
within each sector as far as its counts allow. The reason is statistical power: the population is
heavily skewed toward CDN and offshore hosting (roughly $58\%$ CDN, $26\%$ Abroad, $14\%$ Pakistani
in our $1{,}781$-site candidate pool), so a purely proportional draw would leave on the order of
fourteen Pakistani sites---too few to estimate that class or to detect a difference between classes.
Equalising the per-class sample sizes minimises the variance of a difference of means and therefore
maximises the power of the between-class comparison at the heart of RQ1. Stratifying on both axes at
once lets us compare hosting classes \emph{inside} a sector---government hosted abroad versus
locally, say---without a sector's own mix confounding the effect. Because the hosting marginal is
intentionally non-representative, any population-level figure is recovered by post-stratification,
reweighting the strata back to their true shares; the one sample thus serves both roles, unweighted
for powerful between-class comparison and reweighted for representative population estimates. The
candidate pool from which the sample is frozen is the Tranco popularity ranking filtered to
\texttt{.pk} augmented with curated critical-infrastructure sectors, each site's hosting resolved by
DNS then AS/geo classification; seven sites carry no Tranco rank and are included as sector
representatives, and the frozen list records site, sector, hosting class, host, and rank so the
panel is reproducible and auditable.

\begin{table}[t]
  \centering
  \caption{Genre axis: CISA sectors, drawn proportionally to the resolved \texttt{.pk}
  population. The hosting axis is held at a fixed $40$/$40$/$20$ (Pakistani/CDN/Abroad) across the
  sample for power, and reweighted to the population shares for any population-level estimate.}
  \label{tab:sectors}
  \begin{tabular}{lc}
    \toprule
    CISA sector & Sites \\
    \midrule
    Commercial Facilities            & 60 \\
    Government Services \& Facilities & 11 \\
    Education                        & 10 \\
    Communications                   & 7  \\
    Financial Services               & 4  \\
    Energy                           & 3  \\
    Healthcare \& Public Health      & 3  \\
    Transportation Systems           & 2  \\
    \midrule
    Total                            & 100 \\
    \bottomrule
  \end{tabular}
\end{table}

\subsection{Measurements and indicators}
Against each target we run a \textbf{TCP/80 Paris traceroute} once per hour and an ICMP
\textbf{ping} every $30$~minutes, each fanning out to all discovered probes, for a seven-day window.
TCP/80 is used for the path because ICMP is heavily filtered in Pakistani networks and undercounts
badly---in our first tromboning target only $12$ of $52$ destinations answered ICMP---while Paris
addressing holds the flow tuple constant so that an observed path change is a genuine reroute rather
than load-balancer noise. Ping supplies the end-to-end RTT, and we take the \textbf{minimum of the
$N$ packets} as the level because network noise is one-directional---queuing only ever adds
delay---so the minimum is the cleanest estimate of the true path latency, in contrast to the noisy
single-packet RTT of our earliest runs. From these we derive the four indicators used throughout,
following prior IXP work~\cite{dibartolomeo2015}: RTT (the min-of-$N$ level), \textbf{jitter} as the
round-to-round standard deviation, packet loss as sent minus received, and hop count from the
traceroute. The hourly cadence for the heavier traceroute samples path changes without hammering
the network or the credit budget, and the finer $30$-minute ping cadence gives the RTT, jitter, and
loss series the temporal resolution a diurnal analysis needs. Routing is classified not by geo-IP,
which mislabels the very hops that matter, but by the RTT-physics detector we validated in earlier
work and reuse unchanged here (Sec.~\ref{sec:detector}).

\subsection{The tromboning detector (reused)}
\label{sec:detector}
Because registration country lies exactly where it counts---a ``Canadian'' Shaw address and a
``US'' Cogent address both sit physically in Pakistan at ${\sim}2$\,ms, while a genuinely foreign
exit hop is often unannounced in BGP---we decide tromboning by RTT physics rather than hop country.
A trace is judged to have left Pakistan if any responding hop is foreign with RTT $\geq 40$\,ms, or
there is a $\geq 60$\,ms jump between consecutive hops, or any hop RTT reaches $\geq 70$\,ms; a path
whose maximum RTT stays under $45$\,ms is local. RTTs above $500$\,ms are discarded as queuing or
ICMP-error artefacts, and the Shaw/Cogent addresses are excluded. This is the same detector that, on
our first ISP, cut a naive geo-IP test's $53/53$ false positives to a clean $16/52$ matching manual
inspection; carrying it unchanged into PANEL is deliberate, so that the panel's tromboning verdicts
are directly comparable to the census that preceded it.

\subsection{Operation under the concurrency limit}
\label{sec:ops}
The measurements are registered as server-side \emph{periodic} instructions with a fixed seven-day
window, so collection runs on RIPE's own infrastructure and a local watcher can restart freely
without losing data. The one binding constraint is RIPE Atlas's cap of $100$ simultaneous
measurements \emph{per account}: one hundred sites each needing a traceroute and a ping is two
hundred measurements, twice the cap. The daily credit spend ($\approx 1.3$M against a $10$M ceiling)
and daily result volume ($\approx 108$k against $1$M) are comfortable, so the concurrency cap, not
cost, is what shapes the deployment. We resolve it by splitting the panel across two credit-sharing
accounts \emph{by measurement type}: account~A registers all one hundred traceroutes and account~B
all one hundred pings, each exactly at the cap, both against the identical target list and the same
live-discovered probes, with outputs namespaced so both run on one host and merge at analysis on
probe, target, and timestamp. Because the split is by type rather than by target, the merged dataset
is identical to what a single uncapped account would have produced. A preflight guard refuses to
create measurements if the plan would exceed the cap, so the scheme is safe to launch. This paper's
run was scheduled on 2026-07-11 and auto-stops on 2026-07-18; the distilled per-round series is
harvested every $30$~minutes during the window, and the untouched raw results are archived once at
the end, recoverable at any time from the stored measurement identifiers at no additional cost.

\section{Analysis}
\label{sec:analysis}
\emph{Data collection is in progress; the seven-day window closes on 2026-07-18. This section fixes
the analyses the panel is designed to deliver and reserves their figures, so that the reported
numbers follow directly from the design above once the merged dataset is complete.} \pending

\subsection{Per-ISP quality distributions (RQ1)}
For each ISP we will present the distribution of every indicator---RTT, hop count, jitter, and
packet loss---separately to local, CDN, and offshore targets, as cumulative distribution functions
in the manner of prior IXP work, so that operators can be compared on equal terms rather than by a
single average. The central comparison is between ISPs that peer locally and those that do not: if
local peering helps, its customers' curves to local and CDN targets should sit to the left of, and
be steeper (less variable) than, those of non-peering operators. Because the panel holds targets and
method constant, any separation between the curves is attributable to the ISP, which is the
comparison our earlier snapshots could not support.
% \begin{figure}[t] \includegraphics[width=\columnwidth]{fig_panel_rtt_cdf_by_isp.png}
%   \caption{Per-ISP CDF of min-of-$N$ RTT to local vs CDN vs offshore targets.} \label{fig:panel_cdf} \end{figure}
\emph{Reserved:} per-ISP RTT/jitter/loss/hop CDFs; headline result \XX. \pending

\subsection{Same-ISP consistency (RQ2)}
Using the several probes we hold for the same operator, we will quantify how much the indicators
vary \emph{within} an ISP---across cities and access points---to test whether customers of one ISP
receive similar service or whether quality is a property of the particular PoP. Our census already
hinted that routing can be per-PoP rather than per-ISP; the panel measures the spread directly and
over time, and where the same ISP appears in more than one city we normalise across cities before
comparing, so that a genuine same-ISP inconsistency is not mistaken for a mere distance effect.
\emph{Reserved:} within-ISP variance of each indicator; per-PoP breakdown. \pending

\subsection{Tromboning stability over the week}
Applying the reused detector to every traceroute, we will report, for each PK-hosted site that ever
hairpins, the \emph{fraction of rounds} in which it trombones---turning the previously open question
of intermittency into a measured quantity. This directly addresses the caveat that dogged the census,
where a single snapshot could not tell a stable structural hairpin from a verdict that flickers hour
to hour; seven days of hourly traces let us classify each path as persistently local, persistently
hairpinned, or flapping, and report the census rate as a range rather than a point.
\emph{Reserved:} per-path trombone-fraction distribution; stable-vs-flapping split. \pending

\subsection{Diurnal and weekly structure}
Because the panel spans a full week at fine cadence, we will decompose each indicator into
time-of-day and day-of-week components to test whether the offshore penalty and the trombone rate
carry a diurnal or weekday cycle or are flat. A flat profile would confirm the penalty is
structural---a hosting and peering choice rather than evening congestion---strengthening the case
that the fix is local hosting and active peering, not more bandwidth; a cycle, conversely, would
localise where and when the cost is worst.
\emph{Reserved:} diurnal/weekly decomposition of RTT and trombone rate. \pending

\subsection{A standing baseline for the next fault}
Finally, the seven-day series is itself a normal-day baseline. Our earlier outage study could only
compare a cable fault's peak against its own recovered state, for want of a pre-event reference; by
running continuously against a fixed panel, PANEL holds that reference ready, so that if a routing or
submarine-cable event occurs in the window---or against which a future event is later compared---the
disruption is measurable as a true departure from normal, with the domestic and offshore legs
separated as in the outage study.
\emph{Reserved:} baseline indicator levels; event delta if one occurs. \pending

\section{Conclusion and Future Work}
\label{sec:conclusion}
PANEL is the measurement our earlier work implied but could not itself provide: a single, uniform,
seven-day panel that fixes the vantage set, the protocol, the RTT definition, and the metric set, so
that the per-ISP quality of Pakistani Internet access can be compared fairly and its stability
measured rather than assumed. We have presented its design, the two-axis sampling that makes it both
representative and powerful, and the two-account scheme that runs it within RIPE Atlas's limits; the
seven-day dataset is in collection, and the per-ISP distributions, the same-ISP consistency test, the
tromboning-stability profile, the diurnal decomposition, and the standing baseline of
Section~\ref{sec:analysis} will be reported from the merged results once the window closes. Beyond
this run, the same framework extends naturally to a longer horizon and to repeated tromboning
censuses, and---by holding a normal-day reference continuously---turns the next cable fault from a
missed baseline into a clean before/during/after.

\section*{Acknowledgment}
Supported by the APNIC Foundation and supervised by Dr.\ Saqib Ilyas at LUMS. Measurements use the
RIPE Atlas platform.

\begin{thebibliography}{99}
\bibitem{dibartolomeo2015} M.~Di Bartolomeo, G.~Di Battista, R.~di Lallo, and C.~Squarcella,
``Is it really worth to peer at IXPs? A comparative study,'' in \emph{Proc. IEEE ISCC}, 2015.
\bibitem{gupta2014} A.~Gupta, M.~Calder, N.~Feamster, M.~Chetty, E.~Calandro, and
E.~Katz-Bassett, ``Peering at the Internet's frontier: A first look at ISP interconnectivity in
Africa,'' in \emph{Proc. PAM}, 2014.
\bibitem{sanchez2014} M.~A.~Sanchez, F.~E.~Bustamante, B.~Krishnamurthy, W.~Willinger,
G.~Smaragdakis, and J.~Erman, ``Inter-domain traffic estimation for the outsider,'' in
\emph{Proc. ACM IMC}, 2014.
\bibitem{klick2016} J.~Klick, S.~Lau, M.~W\"ahlisch, and V.~Roth, ``Selective, prefix-based
scanning of the IPv4 address space,'' in \emph{Proc. ACM IMC}, 2016. % verify exact title
\bibitem{ager2012} B.~Ager, N.~Chatzis, A.~Feldmann, N.~Sarrar, S.~Uhlig, and W.~Willinger,
``Anatomy of a large European IXP,'' in \emph{Proc. ACM SIGCOMM}, 2012.
\bibitem{chatzis2013} N.~Chatzis, G.~Smaragdakis, A.~Feldmann, and W.~Willinger, ``There is more
to IXPs than meets the eye,'' \emph{ACM SIGCOMM Comput. Commun. Rev.}, 2013.
\bibitem{ripeatlas} RIPE NCC, ``RIPE Atlas,'' \url{https://atlas.ripe.net/}.
\bibitem{ripestat} RIPE NCC, ``RIPEstat,'' \url{https://stat.ripe.net/}.
\end{thebibliography}

\end{document}
